\documentclass[11pt]{article}
\usepackage{fontspec}
\usepackage[french,english]{babel}
\usepackage{amsmath}
\usepackage{amssymb}
\usepackage{amsthm}
\usepackage{mathrsfs}
\usepackage{tikz-cd}
\usepackage{enumitem}
\usetikzlibrary{babel}

\newcommand{\printedpage}[1]{\marginpar{\footnotesize p.~#1}}
\newcommand{\msp}[1]{\marginpar{\footnotesize [#1]}}
\newcommand{\circled}[1]{\tikz[baseline=(c.base)]\node[draw,circle,inner sep=1pt](c){#1};}
\newcommand{\Cof}{\mathrm{Cof}}
\newcommand{\Fib}{\mathrm{Fib}}
\newcommand{\Fl}{\mathrm{Fl}}
\newcommand{\univ}{\mathrm{univ}}

\begin{document}

\printedpage{9}
% head: A. Grothendieck : Dérivateurs, Ch. XIII, Édition des Universités de Montpellier II et Paris VII

\begin{itemize}[leftmargin=2.5em]
\item[\circled{1}] \textit{Then} $\Cof_W$ \textit{and} $W^{\univ}$ \textit{have the stability properties} a) \textit{and} c) \textit{of proposition} 2 (p.~4), \textit{and also the stability property} b) \textit{under direct factors, provided that} $W$ \textit{itself has it.}

\item[\circled{2}] \textit{Suppose} $\mathcal{M}$ accessible, \textit{so that we have theorem} 1. \textit{Equivalent conditions on} $W$ \textit{:}

\msp{page 15}

\begin{itemize}[leftmargin=2.5em]
\item[a)] \textit{There is a small subset} $C_0$ \textit{of} $\Cof_W$\textit{, and} $TC_0$ \textit{of} $W^{\univ}$\textit{, such that}
\[ TF \overset{\text{def}}{=} C_{0*} \subset W, \qquad (TC_0)_* \overset{\text{def}}{=} F \subset \Fib_W. \]

\item[a$'$)] \textit{There is a quadruple} $(C,TC,F,TF)$ \textit{of subsets of} $\Fl(\mathcal{M})$\textit{, satisfying conditions} M~1 \textit{to} M~4, M$'$, M$''$ \textit{on page} 6 (\footnote[10]{with $C$, $TC$ left closed, $F$, $TF$ right closed.}).

\item[a$''$)] \textit{As in} a$'$)\textit{, assuming} M$'$ \textit{in strengthened form with} M$'_0$ (g), i.e. \textit{we have a left pre-Q-category of models structure (}$TC = C \cap W$\textit{).}

\item[b)] \textit{We have}
\[ (\Cof_W)_* \ \subset\ W, \qquad (W^{\univ})_* \ \subset\ \Fib_W \]
\textit{(thus} $(\Cof_W)_* \subset \Fib_W \cap W$ \textit{since} $(\Cof_W)_* \subset (W^{\univ})_*$ (\footnote[11]{$W^{\univ} \subset \Cof_W$.})\textit{).}

\item[b$'$)] \textit{We have}
\[ (\Cof_W)_* \ \subset\ W_{\univ}, \qquad (W^{\univ})_* \ \subset\ \Fib_W. \]
\end{itemize}

\item[\circled{3}] \textit{For there to be a quadruple which, with} $W$\textit{, defines a right pre-Q-category of models (}$TF = F \cap W$\textit{)} (\footnote[12]{We are still assuming $\mathcal{M}$ accessible in 3° and 4°, with in addition $F$, $TF$ (and $C$, $TC$) right closed (resp. left closed).})\textit{, it is necessary and sufficient that we have} b$'$) \textit{above, as well as}
\[ (W^{\univ})_* \cap W_{\univ} \ \ \text{\itshape is right closed.} \]

\msp{page 16}

\item[\circled{4}] \textit{The preliminary assumptions on} $(\mathcal{M}, W)$\textit{, including accessibility, as well as the equivalent conditions in}\circled{2}\textit{, are stable under passage from} $(\mathcal{M}, W)$ \textit{to} $(\mathcal{M}(I), W(I))$\textit{. But what about the condition in}\circled{3} \textit{? Is it a consequence of the conditions in} \circled{2} \textit{(}i.e. \textit{is it equivalent to them) ?} % [sic] leurs est-elle -> lui est-elle, grammar
\end{itemize}

I am going to sketch the proof of the theorem-scholie, to pin down the technical work that remains to be done.

1°) Is routine verification.

2°) Here are the $\pm$ tautological implication between the conditions under consideration : % [sic] les implication -> les implications, agreement
\[ a'' \Longrightarrow a' \overset{?}{\Longrightarrow} a \Longrightarrow b \Longleftarrow b', \tag*{$(*)$} \]
% foot: Version 25 janvier 2004
\printedpage{10}
% head: A. Grothendieck : Dérivateurs, Ch. XIII, Édition des Universités de Montpellier II et Paris VII
but I see that $a' \Longrightarrow a$ is not at all tautological, and perhaps even false after all~! (\footnote[13]{That is why I modified it, strengthening $a'$) by \texttt{[the condition of note 10]}.}). We would have to assume that in \texttt{[the]} quadruple $(C,F,TC,TF)$, $F$ and $TF$ are right closed, and more precisely, that $F = (TC)_*$, $TF = C_*$, which will be the case e.g. if we have a ``closed'' Q-category of models structure (``closed model category''), in the sense of \textsc{Quillen}. The other three implications in $(*)$ really are tautological.

$a \Longrightarrow a'$. We can in $a$) assume
\[ TC_0 \ \subset\ C_0 \]
(replacing $C_0$ by $C_0 \cup TC_0$ if necessary), then, thanks to
\msp{page 17}
1°) ($W^{\univ}$ and $\Cof_W$ are left closed), setting
\[ \overline{TC_0} = TC, \qquad \overline{C_0} = C, \qquad \text{\texttt{[rather}}\ \ \widetilde{TC_0} = TC, \qquad \widetilde{C_0} = C, \text{\texttt{]}} \]
we will have $TC \subset W^{\univ}$, $C \subset \Cof_W$. Passing from $TC_0$ to $TC$, from $C_0$ to $C$, we have lost only the fact that these subsets of $\Fl(\mathcal{M})$ are small. But by theorem 1, setting
\[ F \ =\ (TC)_*, \qquad TF \ =\ C_*, \]
the pairs of closed subsets (left -- right)
\[ (C,\, TF), \qquad (TC,\, F) \]
satisfy the Quillen factorization condition. We will then have properties M~1 to M~4, plus the inclusions $TC \ \subset\ C \cap W$ (since $TC \ \subset\ W^{\univ} \ \subset\ W$, and $TC \ \subset\ C$ by construction), and $C \subset \Cof_W$ of M$'$ and M$''$. It remains to check $TF \subset F \cap W$, $F \subset \Fib_W$, but $F$ $(= (TC)_*)$ $\ \subset \Fib_W$ by assumption, and $TF$ $(= C_*)$ $\ \subset W$ too~; on the other hand $TF \subset F$ follows from $TC \subset C$.

$a' \Longrightarrow a''$. We thus want, in addition to conditions M~1 to M~4 and M$'$, M$''$, the condition
\[ TC \ =\ C \cap W\ (= C \cap W^{\univ}). \]
The idea is to replace the pair
\msp{page 18}
\[ C,\ TC \qquad (C \subset \Cof_W,\ TC \subset C \cap W), \]
giving rise to
\[ F = (TC)_*, \qquad TF = C_*, \]
by
\[ C,\ C \cap W\ (\supset TC), \]
% foot: Version 25 janvier 2004
\printedpage{11}
% head: A. Grothendieck : Dérivateurs, Ch. XIII, Édition des Universités de Montpellier II et Paris VII
which does not change the pair $(C, TF)$, but does change the pair $(TC, F)$, replaced by
\[ \underbrace{(TC)'}_{\supset TC} \ =\ \underbrace{C \cap W}_{= C\cap W^{\univ}}, \qquad F' \ =\ \underbrace{(TC')_*}_{\subset F} \]
(we have replaced $F$ by a smaller $F'$) (\footnote[14]{NB $TC'$ will still be left closed as an intersection of two left closed subsets.}). The only problem is whether $(TC', F')$ still satisfies the Quillen factorization condition. These are closed subsets (left for $TC'$, right for $F'$) ``dual'' to one another,
\[ TC' = F'^*, \qquad F' = TC'_* \ , \]
but since we do not know whether $TC'$ is of the form $\widetilde{TC'_0}$, with $TC'_0$ small, we are not sure that the conclusion of theorem 1 applies. But if we assume in $a'$) that $C$ is an \textit{accessible} subset of $\Fl(\mathcal{M})$ (\textit{i.e.}, already being left closed, that $C = \widetilde{C_0}$, $C_0$ small),
\msp{page 19}
then I presume that we will be able to prove that the same is true of $C \cap W$ ($W$ itself being accessible, thanks to the condition of stability under small filtering $\varinjlim$), so theorem 1 really does apply to $TC' = C \cap W$. But if we keep $a'$) as it stands, then we would probably have to prove that $a'$) \textit{implies} $a'$) in the stronger form, with $C$ \textit{accessible.} But since we have $a' \Longrightarrow a$, and we have just established that $a \Longrightarrow a'$ in this stronger form, we win (modulo the conjectural fact assumed about the admissible \texttt{[rather accessible]} subsets of $\Fl(\mathcal{M})$).

It thus remains to prove
\[ b \ \Longrightarrow\ a, \qquad b \ \Longrightarrow\ b', \]
and it suffices to prove the first implication, since we know that $a$) $\Longrightarrow$ $a'$), and that $a'$) implies that
\[ \Fib_W \cap W \ =\ W_{\univ} \]
(cor.~1, page 10), so $(\Cof_W)_* \subset W$ and $(W^{\univ})_* \subset \Fib_W$ (\textit{i.e.} $b$)) imply, since $(\Cof_W)_* \subset (W^{\univ})_*$ (because $W^{\univ} \subset \Cof_W$), $(\Cof_W)_* \subset W \cap \Fib_W = W_{\univ}$, which is the part of $b'$) to be established (the other $(W^{\univ})_* \subset \Fib_W$ being tautologically contained in $b$).

\msp{page 20}

Let us now note that $b$) is nothing other than the weakened variant of $a$), when we omit specifying that $C_0$, $TC_0$ must be small. So the only problem, for $b \Longrightarrow a$, is to see this :

\textbf{Conjectural lemma} (\footnote[15]{This lemma looks a bit false, but it should be possible to fix that \dots}).

\begin{itemize}[leftmargin=2.5em]
\item[\circled{1}] \textit{If we have a left closed subset} $\Phi$ \textit{of} $\Fl(\mathcal{M})$\textit{, such that} $\Phi_* \subset \Fib_W$ \textit{(resp.} $\Phi_* \subset W$\textit{), then there is a subset} of small size $\Phi_0 \subset \Phi$\textit{, such that} $\Phi_{0*} \subset \Fib_W$ \textit{(resp.} $\Phi_0 \subset W$\textit{).}

\item[\circled{2}] \textit{If} $C$\textit{,} $C'$ \textit{are two left closed subsets of} $\Fl(\mathcal{M})$\textit{, if} $C$ \textit{and} $C'$ \textit{are accessible, the same is true of} $C \cap C'$\textit{.}
\end{itemize}
% foot: Version 25 janvier 2004

\printedpage{12}
% head: A. Grothendieck : Dérivateurs, Ch. XIII, Édition des Universités de Montpellier II et Paris VII
I believe this really is true under the assumptions made on $\mathcal{M}$, $W$ -- which must imply that in a suitable sense, $W$ is an accessible subset of the set $\Fl(\mathcal{M})$ of arrows of the accessible category $\mathcal{M}$. It should follow from this that a subset of $\mathcal{M}$ constructed explicitly in terms of $W$, such as $\Fib_W$, is itself accessible, and the lemma should come out of that.

Thus, 2° is proved, modulo this accessibility lemma, which I will have to prove $\pm$ at the same time as theorem 1 (which I have only stated here).

\msp{page 21}

Let us prove 3°. First suppose there is a ``closed'' quadruple $(C, TC, F, TF)$, fitting into a ``closed Quillen pre-quintuple'' which is a right Q-category of models, \textit{i.e.}
\[ TF \ =\ F \cap W \ =\ F \cap W_{\univ}. \]
Let us then prove \texttt{[that]}
\[ \underbrace{(W^{\univ})_* \cap W_{\univ}}_{A_W} \qquad \text{is right closed.} \tag*{$(*)$} \]
Indeed, we have
\[ TC \ \subset\ W^{\univ}, \]
hence
\[ (W^{\univ})_* \ \subset\ (TC)_* \ =\ F, \]
so
\[ A = (W^{\univ})_* \cap W_{\univ} \ \subset\ F \cap W_{\univ} \ \overset{\text{(by hyp.)}}{=}\ TF, \]
now $TF$ is right closed and contained in $W_{\univ}$, hence
\[ \overline{A} \ \subset\ TF \ \subset\ W_{\univ}, \]
and since $(W^{\univ})_*$ is right closed, we also have
\[ \overline{A} \ \subset\ (W^{\univ})_*, \]
hence $\overline{A} \subset W_{\univ} \cap (W^{\univ})_* = A$, so $A$ is right closed, q.e.d.

Let us prove the converse, \textit{i.e.} suppose the conditions of 2° are satisfied, \textit{plus} condition $(*)$ above. We must construct the quadruple.

\msp{page 22}

Let us first note, to get our bearings, that the relations $b'$) of 2°), and the transposed relations (taking into account that $\Cof_W$ and $W^{\univ}$ are left closed, by 1°), namely
\[ \left\{ \begin{array}{lllllll}
(\Cof_W)_* & \subset & W_{\univ} & , & (W^{\univ})_* & \subset & \Fib_W \\
(W_{\univ})^* & \subset & \Cof_W & , & (\Fib_W)^* & \subset & W^{\univ},
\end{array} \right. \]
giving (\footnote[16]{taking into account $W^{\univ} \subset \Cof_W$, $W_{\univ} \subset \Fib_W$ and the transposed inclusions.}) the two square diagrams of inclusions
% foot: Version 25 janvier 2004
\printedpage{13}
% head: A. Grothendieck : Dérivateurs, Ch. XIII, Édition des Universités de Montpellier II et Paris VII
\[
\begin{tikzcd}[column sep=small, row sep=small]
 & W^{\univ} \arrow[dr, hook] & \\
(\Fib_W)^* \arrow[ur, hook] \arrow[dr, hook] & & \Cof_W \\
 & (W_{\univ})^* \arrow[ur, hook] &
\end{tikzcd}
\qquad\qquad (\text{\footnotemark[17]})
\tag*{(I)}
\]
\footnotetext[17]{\textbf{NB} We have these squares of inclusions, whenever we have a Quillen system $(W, C, F, TC, TF)$ that is \textit{closed.}}
and
\[
\begin{tikzcd}[column sep=small, row sep=small]
 & (W^{\univ})_* \arrow[dl, hook'] & \\
\Fib_W & & (\Cof_W)_*\ , \arrow[ul, hook'] \arrow[dl, hook'] \\
 & W_{\univ} \arrow[ul, hook', "\alpha"'] &
\end{tikzcd}
\tag*{(II)}
\]
of which the first is the transpose (under $\Phi \longmapsto \Phi^*$) of the second, and where the 8 sets of arrows considered, except for $\Fib_W$ and $W_{\univ}$ (linked by $\alpha$ in square II) are closed (the four sets in I left closed, and the two sets in II not involved in $\alpha$, right closed).

\msp{page 23}

The ``closed Quillen pre-quintuples'', associated with $W$, correspond to the choice of left closed subsets $TC$ and $C$, fitting into diagram (I) as shown in I$'$ below, and give rise by transposition to $F$ and $TF$ as in II$'$ below,
\[
\begin{tikzcd}[column sep=small, row sep=small]
 & & W^{\univ} \arrow[dr, hook] & \\
 & TC \arrow[ur, hook] \arrow[dr, hook] & & \Cof_W \\
(\Fib_W)^* \arrow[ur, hook] \arrow[dr, hook] & & C \arrow[ur, hook] & \\
 & (W_{\univ})^* \arrow[ur, hook] & &
\end{tikzcd}
\tag*{(I$'$)}
\]
\[
\begin{tikzcd}[column sep=small, row sep=small]
 & & (W^{\univ})_* \arrow[dl, hook'] & \\
 & F \arrow[dl, hook'] & & (\Cof_W)_* \arrow[ul, hook'] \arrow[dl, hook'] \\
\Fib_W & & TF \arrow[ul, hook'] \arrow[dl, hook'] & \\
 & W_{\univ}\ , \arrow[ul, hook'] & &
\end{tikzcd}
\tag*{(II$'$)}
\]
the ``catch'' thus being to choose $TC$ and $C$ large enough that we have
\[ \underbrace{(TC)_*}_{= F} \ \subset\ \Fib_W, \qquad \underbrace{C_*}_{= TF} \ \subset\ W_{\univ} \qquad \text{--} \]
% foot: Version 25 janvier 2004

\end{document}
